%% !TEX root = manuscript.tex

\chapter*{Introduction}

Le groupe Angelotti est un acteur de l'aménagement et de la promotion immobilière implanté en
Occitanie et en région PACA. Ces deux métiers n'ont pas la même logique économique : dans
l'aménagement, la valeur naît de la transformation du foncier~; dans la promotion, de la
construction. Ils produisent deux structures de coûts et deux rythmes de commercialisation
différents, et cette frontière ressortira d'elle-même de l'analyse exploratoire du chapitre~3.
Chaque opération étant par ailleurs portée par une société dédiée, le suivi financier est
nativement cloisonné et toute vision transverse doit être reconstruite.

Une opération vit sur un cycle qui se compte en années. Un budget est arrêté à l'engagement,
puis le chantier et la commercialisation le font dériver. Cette dérive est le fonctionnement
normal du métier~; ce qui pose problème, c'est de la mesurer à temps, car sur une marge de
quelques centaines de milliers d'euros, quelques mois de retard suffisent à transformer un
ajustement en perte.

À mon arrivée en deuxième année, le système d'information reposait sur trois outils ---
\textbf{Grimmo}, l'ERP métier historique, pour la gestion des opérations~; \textbf{Vadimm}, le CRM, pour
la relation commerciale~; \textbf{MyReport}, l'outil décisionnel, pour les tableaux de bord --- et sur
une quatrième couche moins visible, \textbf{les tableurs}, qui faisaient le lien partout où les
trois premiers ne se parlaient pas.

Le contexte data de mes missions se résume alors à un paradoxe que j'ai mis plusieurs mois à
formuler : l'entreprise ne manquait pas de données. Elle disposait de l'historique complet de
ses opérations, ventes, désistements et budgets sur plusieurs années. Ce qui manquait, c'était
la possibilité de s'en servir, et trois obstacles se cumulaient. L'\textbf{absence d'identifiant
partagé} d'abord : un acquéreur du CRM et un tiers de l'ERP ne se reconnaissent pas, et les
rapprocher suppose de reconstruire une clé d'identité. L'\textbf{absence de définition partagée}
ensuite : la marge d'une opération peut se calculer sur le budget, sur l'engagé ou sur le
facturé, et tant que la définition n'est écrite nulle part, chaque tableau de bord porte
implicitement la sienne, deux services pouvant afficher deux chiffres différents sans que ni
l'un ni l'autre ait tort. L'\textbf{absence de profondeur temporelle} enfin : les exports du système
de gestion disent où en est un budget aujourd'hui, pas comment il y est arrivé. Ces trois
obstacles ont un point commun : aucun ne relève de la modélisation. Ce sont des problèmes de
structure, et ils se règlent avant tout modèle.

C'est dans ce décor que s'inscrit mon année. Après une première année d'alternance consacrée à
la migration du CRM Vadimm, ma deuxième année de Master MIASHS au sein de la Holding SPA a
coïncidé avec l'étape la plus lourde de la modernisation du système d'information : la refonte
complète de l'ERP métier, passant de Grimmo à la solution SPO éditée par Salvia. Un changement
d'ERP n'est pas un changement d'outil, c'est un changement de modèle de données : les objets ne
portent plus les mêmes noms, ne s'articulent plus de la même façon et n'acceptent plus les
mêmes valeurs. Ce changement d'échelle m'a conduite d'un rôle d'assistance opérationnelle vers
des missions d'ingénierie, centrées sur l'interopérabilité des systèmes.

Rattachée à la Holding SPA, j'ai travaillé au contact direct des services métiers --- direction
financière, service juridique, équipes commerciales --- tout en intervenant sur les outils du
système d'information. Cette position intermédiaire a une conséquence pratique constante : la
formulation du besoin fait partie du travail. Un service ne demande presque jamais un livrable
technique~; il décrit une gêne, un tableau qu'il refait à la main chaque mois, une question à
laquelle il ne sait pas répondre.

Ce mémoire présente les travaux réalisés durant cette année, du premier jour de la migration
jusqu'à la livraison d'une application d'aide à la décision. Il témoigne d'un mouvement en deux
temps : le passage de la Data Science exploratoire au Data Engineering appliqué, puis le retour
à la modélisation une fois les données reconstituées. L'ordre compte, et c'est de lui que je
tire la thèse de ce mémoire, énoncée d'emblée pour que le lecteur puisse en juger au fil des
chapitres : \textbf{la structure de la donnée commande ce que la modélisation peut en tirer}. Ce
n'est pas une position théorique, c'est un constat que l'année a produit deux fois. La première
en creux, lorsque la charge de la migration a repoussé le projet de Machine Learning au second
semestre. La seconde en plein, lorsque ce projet a donné des résultats solides là où les
données avaient une profondeur temporelle, et modestes là où elles n'étaient qu'une
photographie à date.

Quatre chapitres suivent la progression de mes responsabilités. Le chapitre 1 expose le
chantier le plus lourd de l'année, la reprise des données vers le nouvel ERP, d'abord les
opérations puis les tiers. Le chapitre 2 présente les missions filées menées en parallèle sur
les deux semestres : le maintien du système décisionnel et la contribution au déploiement du
Single Sign-On. Le chapitre 3 expose le projet de Data Science conduit au second semestre, le
Copilote Financier. Le chapitre 4, enfin, livre la réflexion que ces missions ont produite sur
le métier, et présente les perspectives et mon projet professionnel.
